% GENERATED FILE. Edit canonical lessons, then run npm run generate:books.
% canonical-source-hash: fnv1a64:8983aceb75082dbf
% canonical-lessons: FR-C110-droite, FR-C110-gauche, FR-C110-monde, FR-C110-voyage, FR-C110-vacances

\chapter{Right, Left, World, Journey, Holidays}
\label{ch:fr-right-left-world-journey-holidays}
\hlchaptermodality{110}

\begin{chapteropening}
\textbf{By the end of this chapter:} \emph{I can name the right, the left, the world, a journey, holidays.}

Complete the last lesson of chapter 110: I can name the right, the left, the world, a journey, holidays.
\end{chapteropening}

\section[la droite]{la droite — the right}

\label{lesson:FR-C110-droite}



\begin{quote}
Before the new one: say the French for the west, then the French for a letterbox.
\end{quote}

\subsection*{You'll want to know: la droite}
\textbf{la droite} — \textquotedblleft{}the right\textquotedblright{}.

\begin{cousinweb}
\textbf{droite} is the feminine of the word for \textquotedblleft{}right, straight\textquotedblright{}, $\leftarrow$ Latin \emph{directus}, \textquotedblleft{}straight\textquotedblright{} (Lewis \& Short) — inherited: \emph{-ct-} became \emph{-it}.
\end{cousinweb}

\begin{grammarlens}[title={What you've built}]
One more word for this chapter.
\end{grammarlens}

\subsection*{Your turn}
\begin{itemize}
\raggedright
  \item \emph{Say it:} \emph{la droite}
  \item \emph{Say it:} \emph{la droite}, once more
  \item \emph{Say it:} say \emph{la droite}
  \item \emph{Recall:} say the French for the west, then the French for a letterbox, then say \emph{la droite} again
\end{itemize}

\begin{tcolorbox}[breakable,colback=teal!4,colframe=teal!35!black,title={Before you move on}]
What does la droite mean? (\textquotedblleft{}The right\textquotedblright{}.) Say it once more.
\end{tcolorbox}
\section[la gauche]{la gauche — the left}

\label{lesson:FR-C110-gauche}



\begin{quote}
Before the new one: say the French for a letterbox, then the French for the right.
\end{quote}

\subsection*{You'll want to know: la gauche}
\textbf{la gauche} — \textquotedblleft{}the left\textquotedblright{}.

\begin{cousinweb}
\textbf{gauche} — formed from the Old French verb \emph{gauchir}, \textquotedblleft{}to veer, to warp\textquotedblright{} (EtymDB 2.1), usually traced to Frankish \emph{*walkan}, \textquotedblleft{}to trample\textquotedblright{}. The left was the clumsy hand.
\end{cousinweb}

\begin{grammarlens}[title={What you've built}]
One more word for this chapter.
\end{grammarlens}

\subsection*{Your turn}
\begin{itemize}
\raggedright
  \item \emph{Say it:} \emph{la gauche}
  \item \emph{Say it:} \emph{la gauche}, once more
  \item \emph{Say it:} say \emph{la gauche}
  \item \emph{Recall:} say the French for a letterbox, then the French for the right, then say \emph{la gauche} again
\end{itemize}

\begin{tcolorbox}[breakable,colback=teal!4,colframe=teal!35!black,title={Before you move on}]
What does la gauche mean? (\textquotedblleft{}The left\textquotedblright{}.) Say it once more.
\end{tcolorbox}
\section[le monde]{le monde — the world}

\label{lesson:FR-C110-monde}



\begin{quote}
Before the new one: say the French for the right, then the French for the left.
\end{quote}

\subsection*{You'll want to know: le monde}
\textbf{le monde} — \textquotedblleft{}the world\textquotedblright{}.

\begin{cousinweb}
\textbf{monde} $\leftarrow$ Latin \emph{mundus}, \textquotedblleft{}the world\textquotedblright{} — inherited. (EtymDB 2.1)
\end{cousinweb}

\begin{grammarlens}[title={What you've built}]
One more word for this chapter.
\end{grammarlens}

\subsection*{Your turn}
\begin{itemize}
\raggedright
  \item \emph{Say it:} \emph{le monde}
  \item \emph{Say it:} \emph{le monde}, once more
  \item \emph{Say it:} say \emph{le monde}
  \item \emph{Recall:} say the French for the right, then the French for the left, then say \emph{le monde} again
\end{itemize}

\begin{tcolorbox}[breakable,colback=teal!4,colframe=teal!35!black,title={Before you move on}]
What does le monde mean? (\textquotedblleft{}The world\textquotedblright{}.) Say it once more.
\end{tcolorbox}
\section[le voyage]{le voyage — a journey, a trip}

\label{lesson:FR-C110-voyage}



\begin{quote}
Before the new one: say the French for the left, then the French for the world.
\end{quote}

\subsection*{You'll want to know: le voyage}
\textbf{le voyage} — \textquotedblleft{}a journey, a trip\textquotedblright{}.

\begin{cousinweb}
\textbf{voyage} $\leftarrow$ Old French \emph{voiage} $\leftarrow$ Latin \emph{viaticum}, \textquotedblleft{}money and food for a journey\textquotedblright{}, from \emph{via}, \textquotedblleft{}a road\textquotedblright{} — inherited. (EtymDB 2.1)
\end{cousinweb}

\begin{grammarlens}[title={What you've built}]
One more word for this chapter.
\end{grammarlens}

\subsection*{Your turn}
\begin{itemize}
\raggedright
  \item \emph{Say it:} \emph{le voyage}
  \item \emph{Say it:} \emph{le voyage}, once more
  \item \emph{Say it:} say \emph{le voyage}
  \item \emph{Recall:} say the French for the left, then the French for the world, then say \emph{le voyage} again
\end{itemize}

\begin{tcolorbox}[breakable,colback=teal!4,colframe=teal!35!black,title={Before you move on}]
What does le voyage mean? (\textquotedblleft{}A journey, a trip\textquotedblright{}.) Say it once more.
\end{tcolorbox}
\section[les vacances]{les vacances — holidays}

\label{lesson:FR-C110-vacances}



\begin{quote}
Before the new one: say the French for the world, then the French for a journey.

The holiday month: say \emph{en août}.
\end{quote}

\subsection*{You'll want to know: les vacances}
\textbf{les vacances} — \textquotedblleft{}holidays\textquotedblright{}. Always plural in this sense.

\begin{cousinweb}
\textbf{vacances} $\leftarrow$ Latin \emph{vacantia}, \textquotedblleft{}empty, free times\textquotedblright{}, from \emph{vacare}, \textquotedblleft{}to be empty, to be free\textquotedblright{} (Lewis \& Short) — a learned borrowing.
\end{cousinweb}

\begin{grammarlens}[title={What you've built}]
One more word for this chapter.
\end{grammarlens}

\subsection*{Your turn}
\begin{itemize}
\raggedright
  \item \emph{Say it:} \emph{les vacances}
  \item \emph{Say it:} \emph{les vacances}, once more
  \item \emph{Say it:} say \emph{les vacances}
  \item \emph{Recall:} say the French for the world, then the French for a journey, then say \emph{les vacances} again
\end{itemize}

\begin{tcolorbox}[breakable,colback=teal!4,colframe=teal!35!black,title={Before you move on}]
What does les vacances mean? (\textquotedblleft{}Holidays\textquotedblright{}.) Say it once more.
\end{tcolorbox}
